\documentclass[slideColor]{prosper}
\usepackage{amsmath,amsfonts,amsthm}
%\usepackage{epsfig}
%\include{Qcircuit}


\newcommand{\bbC}{{\mathbb{C}}}
\newcommand{\C}{{\mathbb{C}}}
\newcommand{\F}{{\mathbb{F}}}
\newcommand{\N}{{\mathbb{N}}}
\newcommand{\Q}{{\mathbb{Q}}}
\renewcommand{\R}{{\mathbb{R}}}
\newcommand{\Z}{{\mathbb{Z}}}

\newcommand{\cA}{{\mathcal{A}}}
\newcommand{\cD}{{\mathcal{D}}}
\newcommand{\cH}{{\mathcal{H}}}
\newcommand{\cN}{{\mathcal{N}}}
\newcommand{\cP}{{\mathcal{P}}}
\newcommand{\cQ}{{\mathcal{Q}}}
\newcommand{\cS}{{\mathcal{S}}}
\newcommand{\cU}{{\mathcal{U}}}
\newcommand{\cV}{{\mathcal{V}}}

\newcommand{\tr}{\mathop{\mathrm{tr}}}
\newcommand{\rank}{\mathop{\mathrm{rank}}}
\newcommand{\poly}{\mathop{\mathrm{poly}}}
\newcommand{\E}{\mathop{\mbox{$\mathbf{E}$}}}
\newcommand{\schur}{\mathop{\mathrm{Schur}}}
\newcommand{\schursmall}[2]{{\mathrm{S}}_{#1,#2}}
\newcommand{\planch}{\mathop{\mathrm{Planch}}}
\newcommand{\planchsmall}[1]{{\mathrm{P}}_{#1}}
\renewcommand{\d}{{\mathrm{d}}}

\newcommand{\HSP}{\textsc{hsp}\xspace}

\renewcommand{\>}{\rangle}
\newcommand{\<}{\langle}
\newcommand{\ket}[1]{|#1\rangle}
\newcommand{\bra}[1]{\langle #1|}
\newcommand{\scalar}[2]{\langle #1|#2\rangle}
\newcommand{\proj}[1]{\left|#1\right\>\!\left\<#1\right|}
\newcommand{\norm}[1]{\|#1\|}
\newcommand{\ot}{\otimes}
\newcommand{\eps}{\epsilon}
\newcommand{\ra}{\rightarrow}

\renewcommand{\setminus}{\mathbin{-}}
\newcommand{\partitionof}{\mathbin{\vdash}}


\newcommand{\I}{\mathrm{I}}
\renewcommand{\H}{\mathrm{H}}
\renewcommand{\S}{\mathrm{S}}
\newcommand{\T}{\mathrm{T}}
\newcommand{\G}{\mathrm{G}}
%

\newcommand{\bI}{\mathbb{I}}
\newcommand{\bH}{\mathbb{H}}
\newcommand{\bS}{\mathbb{S}}
\newcommand{\bT}{\mathbb{T}}




%%%%%%%%%%%%%%%%%%%%%%%%%%%%%%%%%%%%%%%%%%%%%%%%%%%%%%%%%%%%

\title{A Simple PromiseBQP-complete\\ Matrix Problem}
\author{Pawel M. Wocjan\\[1ex]School of Electrical Engineering and Computer Science\\ University of Central Florida\\Orlando}
\email{wocjan@eecs.ucf.edu}

\begin{document}

\maketitle

%%%%%%%%%%%%%%%%%%%%%%%%%%%%%%%%%%%%%%%%%%%%%%%%%%%%%%%%%%%%%%%%%%%%%%%%

\begin{slide}{Joint work with}
\begin{center}
Dominik Janzing
\end{center}
\end{slide}

%%%%%%%%%%%%%%%%%%%%%%%%%%%%%%%%%%%%%%%%%%%%%%%%%%%%%%%%%%%%%%%%%%%%%%%%%

\begin{slide}{Challenges for Quantum Computing}

\begin{itemize}
\item find new efficient quantum algorithms 

(provide more evidence that BQP is strictly greater than BPP)

\medskip
\item	understand the computational power and limitations of quantum computers in comparison to classical computers
\end{itemize}

\end{slide}

%%%%%%%%%%%%%%%%%%%%%%%%%%%%%%%%%%%%%%%%%%%%%%%%%%%%%%%%%%%%%%%%%%%%%%%%%

\begin{slide}{Efficient Simulation of Quantum Models}

\medskip
\begin{itemize}
\item explore what quantum models can be simulated efficiently on a classical computers; see for example:

\medskip
\begin{itemize}
\item {\em Match gates and classical simulation of quantum circuits}, Jozsa and Miyake

\medskip
\item Knill-Gottesman Theorem; {\em Generalised Clifford groups and simulation of associated quantum circuits}, Clark, Jozsa, Linden

\medskip
\item Simulating quantum circuits using the formalism of matrix product states and the formalism of contracting tensor networks;
generalizations in {\em On the simulation of quantum circuit}, Jozsa


\medskip
\item ...
\end{itemize}

\end{itemize}

\end{slide}

%%%%%%%%%%%%%%%%%%%%%%%%%%%%%%%%%%%%%%%%%%%%%%%%%%%%%%%%%%%%%%%%%%%%%%%%%

\begin{slide}{Charaterizing BQP}

\begin{itemize}
\item approximately evaluating the Jones polynomial of plat closures of braids at roots of unity, Aharonov, Jones, Landau

\medskip
\item approximating the Potts model at some points of the Tutte plane, Aharonov, Arad, Eban, Landau

\medskip
\item approximating partition functions of lattice models in certain complex-parameter regimes, Van den Nest, D{\"u}r, Raussendorf, Briegel

\medskip
\item approximately evaluating quadratically signed weight enumerators, Knill, Laflamme

\medskip
\item ...

\end{itemize}

\end{slide}

%%%%%%%%%%%%%%%%%%%%%%%%%%%%%%%%%%%%%%%%%%%%%%%%%%%%%%%%%%%%%%%%%%%%%%%%%

\begin{slide}{Language / Decision Problem}

\begin{itemize}
\item a language $L$ is a subset of $\{0,1\}^*$

\medskip
\item the corresponding decision problem (``language recognition'') consists in determining if

\medskip
\begin{itemize}
\item either $x\in L$ 

\medskip
\item or $x\not\in L$ 
\end{itemize}

\medskip
for given input $x\in\{0,1\}^*$
\end{itemize}
\end{slide}

%%%%%%%%%%%%%%%%%%%%%%%%%%%%%%%%%%%%%%%%%%%%%%%%%%%%%%%%%%%%%%%%%%%%%%%%%

\begin{slide}{Complexity Class BQP}

\begin{itemize}
\item BQP is the class of decision problems that can be solved by a {\em uniform} family of quantum circuits $\{U_n\}$ in the following sense:

\bigskip
let $x\in\{0,1\}^n$ $\Rightarrow$
the application of $U_n$ to $|x,0\>$ yields
\[
U_n |x,0\> = 
\alpha_{x,\textcolor{red}{0}}|\textcolor{red}{0}\> \otimes |\psi_{x,0}\> + 
\alpha_{x,\textcolor{blue}{1}}|\textcolor{blue}{1}\> \otimes |\psi_{x,1}\>
\]
such that 

\medskip
\begin{itemize}
\item $|\alpha_{x,1}|^2\ge 2/3$ if $x\in L$
\item $|\alpha_{x,1}|^2\le 1/3$ if $x\not\in L$
\end{itemize}
\end{itemize}
\end{slide}

%%%%%%%%%%%%%%%%%%%%%%%%%%%%%%%%%%%%%%%%%%%%%%%%%%%%%%%%%%%%%%%%%%%%%%%%%

\begin{slide}{Promise Problem}

\begin{itemize}
\item a promise problem $\Pi$ is a pair $(\Pi_{\rm YES},\Pi_{\rm NO})$ of non-intersecting sets: 

\medskip
$\Pi_{\rm YES},\Pi_{\rm NO}\subseteq\{0,1\}^*$ and $\Pi_{\rm YES}\cap\Pi_{\rm NO}=\emptyset$

\bigskip
\item
the set $\Pi_{\rm YES}\cup\Pi_{\rm NO}$ is called the {\em promise}

\bigskip
\item
standard ``language recognition'' problems can be cast as the special case:

\medskip
$\Pi_{\rm YES}\cup\Pi_{\rm NO}=\{0,1\}^*$ ({\em trivial promise})
\end{itemize}
\end{slide}

%%%%%%%%%%%%%%%%%%%%%%%%%%%%%%%%%%%%%%%%%%%%%%%%%%%%%%%%%%%%%%%%%%%%%%%%%

\begin{slide}{Complexity Class PromiseBQP}
\begin{itemize}
\item PromiseBQP is the class of promise problems $(\Pi_{\rm YES},\Pi_{\rm NO})$ that can be solved by a {\em uniform} family of quantum circuits $\{U_n\}$ in the following sense:

\bigskip
\item let $x\in\{0,1\}^n$ $\Rightarrow$ the application of $U_n$ to $|x,0\>$ yields
\[
U_n |x,0\> = \alpha_{x,0}|0\> \otimes |\psi_{x,0}\> + \alpha_{x,1}|1\> \otimes |\psi_{x,1}\>
\]

such that 

\medskip
\begin{itemize}
\item $|\alpha_{x,1}|^2\ge 2/3$ if $x\in\Pi_{\rm YES}$
\item $|\alpha_{x,1}|^2\le 1/3$ if $x\in\Pi_{\rm NO}$
\end{itemize}

\medskip
\item
nothing is required with respect to $x\not\in\Pi_{\rm YES}\cup \Pi_{\rm NO}$
\end{itemize}
\end{slide}

%%%%%%%%%%%%%%%%%%%%%%%%%%%%%%%%%%%%%%%%%%%%%%%%%%%%%%%%%%%%%%%%%%%%%%%%%

\begin{slide}{Complete Problems}
\begin{itemize}
\item no complete problems are known for complexity classes such as BPP, BQP, and MA

\medskip
\item complete problems are only known for their promise versions 
\end{itemize}
\end{slide}

%%%%%%%%%%%%%%%%%%%%%%%%%%%%%%%%%%%%%%%%%%%%%%%%%%%%%%%%%%%%%%%%%%%%%%%%

\begin{slide}{Diagonal Entry Estimation I}
\begin{itemize}
\item decide if

\medskip
\begin{itemize}
\item $(A^m)_{jj} \ge g + \epsilon b^m$ 

\medskip
\item $(A^m)_{jj} \le g - \epsilon b^m$
\end{itemize}

\medskip
\item $A$ is a sparse symmetric matrix of size $N\times N$ with $\|A\| \le b$ 

\medskip 
\item $j\in [1,\ldots,N]$

\medskip
\item $m={\rm polylog}(N)$

\medskip
\item 
$\epsilon=1/{\rm polylog}(N)$

\medskip
\item $g\in[-b^m,b^m]$
\end{itemize}
\end{slide}

%%%%%%%%%%%%%%%%%%%%%%%%%%%%%%%%%%%%%%%%%%%%%%%%%%%%%%%%%%%%%%%%%%%%%%%%

\begin{slide}{Diagonal Entry Estimation II}
\begin{itemize}
\item the task is to decide if the tuple $(A,b,m,j,\epsilon,g)$ is an element of $\Pi_{\rm YES}$ or of $\Pi_{\rm NO}$

\medskip
\item the instances are defined as follows:

\medskip
\begin{itemize}
\item $(A,b,m,j,\epsilon,g)$ is a YES instance iff
\[
(A^m)_{jj} \ge g + \epsilon b^m
\]

\medskip
\item $(A,b,m,j,\epsilon,g)$ is a NO instance iff
\[
(A^m)_{jj} \le g - \epsilon b^m
\]

\medskip
\item the promise is that only tuples in $\Pi=\Pi_{\rm YES} \cup \Pi_{\rm NO}$ are considered
\end{itemize}
\end{itemize}
\end{slide}

%%%%%%%%%%%%%%%%%%%%%%%%%%%%%%%%%%%%%%%%%%%%%%%%%%%%%%%%%%%%%%%%%%%%%%%%

\begin{slide}{DEE is PromiseBQP-complete}

\begin{itemize}
\item Diagonal Entry Estimation is PromiseBQP-complete

\medskip
\item we prove this by showing that

\medskip
\begin{itemize}
\item DEE is in PromiseBQP $\Leftrightarrow$

\medskip
DEE can be solved efficiently on a quantum computer 

\medskip
\item DEE is PromiseBQP-hard $\Leftrightarrow$

\medskip
quantum circuits can be encoded in such a way that we can decide whether they accept an input or not by looking at the corresponding diagonal entry
\end{itemize}
\end{itemize}

\end{slide}

%%%%%%%%%%%%%%%%%%%%%%%%%%%%%%%%%%%%%%%%%%%%%%%%%%%%%%%%%%%%%%%%%%%%%%%%

\begin{slide}{Spectral Measure I}
\begin{itemize}
\item let 
\[
A = \sum_\lambda \lambda Q_\lambda
\]
be the \textcolor{blue}{spectral decomposition of $A$} and $|\psi\>\in\C^N$ some unit vector

\medskip
\item the \textcolor{red}{spectral measure induced by $A$ and $|\psi\>$} is a probability distribution on the spectrum of $A$ such that 
\[
\Pr(\lambda) = \|Q_\lambda |\psi\> \|^2
\]

\medskip
\item we refer to the process of sampling from the spectral measure as \textcolor{green}{measuring the observable $A$ in the state $|\psi\>$}

\end{itemize}
\end{slide}

%%%%%%%%%%%%%%%%%%%%%%%%%%%%%%%%%%%%%%%%%%%%%%%%%%%%%%%%%%%%%%%%%%%%%%%%

\begin{slide}{Spectral Measure II}
\begin{itemize}
\item the expected value of $A^m$ in the state $|\psi\>$ is given by
\[
\<\psi|A^m|\psi\> = \sum_\lambda \lambda^m \<\psi |Q_\lambda|\psi\>
\]
and is equal to the $m$th statistical moment of the spectral measure

\medskip
\item we have $(A^m)_{jj} = \<j|A^m|j\>$

\bigskip
$\Rightarrow$ find quantum algorithm for estimating statistical moments of spectral measures to show that DEE is in PromiseBQP
\end{itemize}
\end{slide}

%%%%%%%%%%%%%%%%%%%%%%%%%%%%%%%%%%%%%%%%%%%%%%%%%%%%%%%%%%%%%%%%%%%%%%%%

\begin{slide}{Outline of Quantum Algorithm}
\begin{itemize}
\item we measure the observable $A$ in the state $j$

\medskip
\item we obtain the eigenvalue $\lambda$ and compute $\lambda^m$ 

\medskip
\item we repeat this process and compute the average value of these powers
\end{itemize}
\end{slide}

%%%%%%%%%%%%%%%%%%%%%%%%%%%%%%%%%%%%%%%%%%%%%%%%%%%%%%%%%%%%%%%%%%%%%%%%

\begin{slide}{Measuring the Observable $A$}

\begin{itemize}
\item we can measure the observable $A$ by

\medskip
\begin{itemize}
\item considering $A$ as a Hamiltonian and simulating the corresponding time dynamics
\[
U = \exp(i A / \|A\|)
\]

\medskip
\item applying the phase estimation algorithm to $U$
\end{itemize}
\end{itemize}

\end{slide}

%%%%%%%%%%%%%%%%%%%%%%%%%%%%%%%%%%%%%%%%%%%%%%%%%%%%%%%%%%%%%%%%%%%%%%%%

\begin{slide}{Analysis of Error Sources}

\begin{itemize}
\item this quantum algorithm works because the following errors can be made sufficiently small with polynomial resources only

\medskip
\begin{itemize}
\item errors due to imperfections (limited resolution and success probability) of phase estimation

\medskip
\item statistical errors due to estimation of the expected value from the empirical mean

\medskip
\item errors due to imperfect simulation of the Hamiltonian time evolution
\end{itemize}
\end{itemize}
\end{slide}

%%%%%%%%%%%%%%%%%%%%%%%%%%%%%%%%%%%%%%%%%%%%%%%%%%%%%%%%%%%%%%%%%%%%%%%%

\begin{slide}{DEE is PromiseBQP-hard}

\end{slide}

%%%%%%%%%%%%%%%%%%%%%%%%%%%%%%%%%%%%%%%%%%%%%%%%%%%%%%%%%%%%%%%%%%%%%%%%

\begin{slide}{Equivalent Definition of PromiseBQP}

\begin{itemize}
\item let $U$ be a BQP quantum circuit and let $x$ be the input

\medskip
\item let $p_0$ and $p_1$ denote the probability of finding the output qubit in state $|0\>$ and $|1\>$, respectively

\medskip
\item elementary calculations show that
\[
\<x,0|U^\dagger\sigma_z^{\textrm{(ouput)}} U|x,0\>=p_0 - p_1
\]

\medskip
\item the sign of this quantity indicates if $x\in\Pi_{\rm YES}$ or $x\in\Pi_{\rm NO}$

\bigskip
$p_0 - p_1 \le -1/3$ if $x\in\Pi_{\textrm{YES}}$ and $\ge 1/3$ if $x\in\Pi_{\textrm{NO}}$
\end{itemize}

\end{slide}

%%%%%%%%%%%%%%%%%%%%%%%%%%%%%%%%%%%%%%%%%%%%%%%%%%%%%%%%%%%%%%%%%%%%%%%%

\begin{slide}{Forward-time Operator}

\begin{itemize}
\item let $V = U^\dagger \sigma_z^{(\textrm{output})} U$ and
\[
V = V_{M-1} V_{M-2} \cdots V_0
\]
the decomposition of $V$ into elementary gates ($M$ is odd)

\medskip
\item the \textcolor{blue}{forward-time operator $F$} is
\[
F = \sum_{\ell=0}^{M-1} |\ell+1\>\<\ell|_{\mathrm{clock}} \otimes V_\ell
\]

\medskip
\item we have $F^{2M}=I$ because $V^2=I$
\end{itemize}

\end{slide}

%%%%%%%%%%%%%%%%%%%%%%%%%%%%%%%%%%%%%%%%%%%%%%%%%%%%%%%%%%%%%%%%%%%%%%

\begin{slide}{Orbit}

\begin{itemize}
\item we restrict our attention to the orbit
\[
\left\{
F^\ell |0\>_{\mathrm{clock}} \otimes |x,0\> \, : \, \ell\in\N
\right\}
\]

\end{itemize}

\end{slide}

%%%%%%%%%%%%%%%%%%%%%%%%%%%%%%%%%%%%%%%%%%%%%%%%%%%%%%%%%%%%%%%%%%%%%%

\begin{slide}{Action of the Forward-time Operator}

\begin{itemize}
\item the input $x$ is fully rejected ($p_0=1$) $\Rightarrow$

\begin{itemize}
\item the orbit is $M$-periodic 

\medskip
\item the action of $F$ is isomorphic to the action of a cyclic shift in $M$ dimensions
\[
|\ell\> \mapsto |(\ell + 1) \mod M \>
\]
where $F^\ell |0\>\otimes |x,0\>$ corresponds to $|\ell\>$

\medskip
\item the spectral measure $R^{(0)}$ induced by $F$ and $|0\>\otimes|x,0\>$ is 
the uniform distribution on the $M$th roots of unity
\[
\Pr\big(\exp(i\pi 2\ell/M)\big) = \frac{1}{M} 
\]

\end{itemize}
\end{itemize}

\end{slide}

%%%%%%%%%%%%%%%%%%%%%%%%%%%%%%%%%%%%%%%%%%%%%%%%%%%%%%%%%%%%%%%%%%%%%%

%%%% YES CASE %%%%

\begin{slide}{Action of the Forward-time Operator}

\begin{itemize}
\item the input $x$ is fully accepted ($p_1=1$) $\Rightarrow$

\medskip
\begin{itemize}
\item the action of $F$ is isomorphic to the action of a cyclic shift in $M$ dimensions with an \textcolor{red}{additional phase $-1$}
\begin{eqnarray*}
|\ell\> & \mapsto & \phantom{-}|\ell + 1\> \mbox{ for $\ell=0,\ldots,M-2$} \\
|M-1\>  & \mapsto & \textcolor{red}{-} |0\>
\end{eqnarray*}

\medskip
\item the spectral measure $R^{(1)}$ induced by $F$ and $|0\>\otimes|x,0\>$ is 
\[
\Pr\big(\exp(i\pi (2\ell \textcolor{red}{+ 1}) / 2M)\big) = \frac{1}{M} 
\]

\end{itemize}
\end{itemize}

\end{slide}

%%%%%%%%%%%%%%%%%%%%%%%%%%%%%%%%%%%%%%%%%%%%%%%%%%%%%%%%%%%%%%%%%%%%%%

\begin{slide}{Spectral Measure}

\begin{itemize}
\item in the general case, the spectral measure induced by $F$ and $|0\>\otimes |x,0\>$ is a convex
combination
\[
R = p_0 R^{(0)} + p_1 R^{(1)}
\]
\end{itemize}

\end{slide}

%%%%%%%%%%%%%%%%%%%%%%%%%%%%%%%%%%%%%%%%%%%%%%%%%%%%%%%%%%%%%%%%%%%%%%

\begin{slide}{Sparse Matrix}

\begin{itemize}
\item we define the self-adjoint sparse matrix
\[
A = \frac{1}{2} (F + F^\dagger)
\]

\medskip
\item
the spectral measure induced by $A$ and $|0\>\otimes |x,0\>$ is given by
\[
P = p_0 P^{(0)} + p_1 P^{(1)}
\]

\medskip
where the supports of $P^{(0)}$ and $P^{(1)}$ are given by the real parts of $R^{(0)}$ and $R^{(1)}$ (cosine values)
\end{itemize}
\end{slide}

%%%%%%%%%%%%%%%%%%%%%%%%%%%%%%%%%%%%%%%%%%%%%%%%%%%%%%%%%%%%%%%%%%%%%%

\begin{slide}{Spectral Measures}

\begin{itemize}
\item $p_0=1$ $\Rightarrow$ the spectral measure $P^{(0)}$ has the form

\[
\begin{array}{ccc|ccc}
&&& \multicolumn{3}{c}{\mbox{for $\ell=1,\ldots,M-1$}}                                \\ \hline 
\lambda^{(0)}           & = & +1   & \lambda^{(0)}_\ell      & = & +\cos(2\pi \ell / M) \\
\Pr(\lambda^{(0)}_0)    & = & 1/M & \Pr(\lambda^{(0)}_\ell) & = & 2/M                 \\
\end{array}
\]

\medskip
\item $p_1=1$ $\Rightarrow$ the spectral $P^{(1)}$ has the form

\[
\begin{array}{ccc|ccc}
&&& \multicolumn{3}{c}{\mbox{for $\ell=1,\ldots,M-1$}}                                \\ \hline 
\lambda^{(1)}           & = & \textcolor{red}{-}1   & \lambda^{(1)}_\ell      & = & \textcolor{red}{-}\cos(2\pi \ell / M) \\
\Pr(\lambda^{(1)}_0)    & = & 1/M & \Pr(\lambda^{(1)}_\ell) & = & 2/M                 \\
\end{array}
\]
\end{itemize}

\end{slide}

%%%%%%%%%%%%%%%%%%%%%%%%%%%%%%%%%%%%%%%%%%%%%%%%%%%%%%%%%%%%%%%%%%%%%%

\begin{slide}{Statistical Moments}

\begin{itemize}
\item identify $j$ with the label of the initial state $|0\>\otimes |x,0\>$

\medskip
\item the $m$th statistical moment is equal to 
\[
(A^m)_{jj} = (p_0 - p_1) \, \cdot \, \left(
1 + \sum_{\ell=1}^{M-1} \cos^{m}(2\pi \ell/M)
\right)
\]

\medskip 
\item for $m=M^3=\mathrm{poly}(\mbox{\# quantum gates})$, we have
\[
(A^m)_{jj} \approx p_0 - p_1 
\]

\end{itemize}

\end{slide}

%%%%%%%%%%%%%%%%%%%%%%%%%%%%%%%%%%%%%%%%%%%%%%%%%%%%%%%%%%%%%%%%%%%%%%

\begin{slide}{Strengthening of the Result}

\begin{itemize}
\item the DEE problem for sparse matrices $A$ with entries restricted to $0,\pm 1$ remains PromiseBQP-hard

\medskip
\item the problem of estimating off-diagonal terms is also in PromiseBQP

\medskip
\item if $A$ is restricted to have only $0,1$ as entries, then the problem of estimating differences
\[
(A^m)_{ii} - (A^m)_{ij}
\]
is PromiseBQP-complete

\medskip
\item the adjacency matrix $A$ can be obtained from simple string rewriting rules (more explicit than saying that $A$ is sparse)
\end{itemize}

\end{slide}

%%%%%%%%%%%%%%%%%%%%%%%%%%%%%%%%%%%%%%%%%%%%%%%%%%%%%%%%%%%%%%%%%%%%%%

\begin{slide}{Connection to Quantum Algorithms}
quantum algorithms that 

\begin{itemize}
\item implement 
\[
\exp(-i A t)
\]

\medskip
\item and compute
\[
\<\psi|f(A)|\psi\>
\]

\medskip
for sparse $A$ and smooth functions $f$

\medskip
\begin{itemize}
\item {\em Quantum algorithm for solving linear systems of equations}, Harrow, Hassidim, Lloyd

\medskip
\item {\em A quantum algorithm to solve nonlinear differential equations}, Leyton, Osborne

\medskip
\item
\end{itemize}
\end{itemize}
\end{slide}

%%%%%%%%%%%%%%%%%%%%%%%%%%%%%%%%%%%%%%%%%%%%%%%%%%%%%%%%%%%%%%%%%%%%%%

\begin{slide}{Conclusions}
\begin{itemize}
\item a purely classical (quantum free) characterization of BQP in terms of a simple matrix problem

\medskip
\item quantum computers are better at counting and comparing the number of solutions of combinatorial problems (provided that BQP is strictly larger BQP)

\medskip
\item $\Rightarrow$ can we find new quantum algorithms based on this observation?
\end{itemize}
\end{slide}

%%%%%%%%%%%%%%%%%%%%%%%%%%%%%%%%%%%%%%%%%%%%%%%%%%%%%%%%%%%%%%%%%%%%%%

\end{document}

%%%%%%%%%%%%%%%%%%%%%%%
%%%%%%%%%%%%%%%%%%%%%%%
%%%%%%%%%%%%%%%%%%%%%%%
%%%%%%%%%%%%%%%%%%%%%%%
%%%%%%%%%%%%%%%%%%%%%%%
%%%%%%%%%%%%%%%%%%%%%%%
%%%%%%%%%%%%%%%%%%%%%%%
%%%%%%%%%%%%%%%%%%%%%%%
%%%%%%%%%%%%%%%%%%%%%%%
%%%%%%%%%%%%%%%%%%%%%%%
%%%%%%%%%%%%%%%%%%%%%%%
%%%%%%%%%%%%%%%%%%%%%%%
%%%%%%%%%%%%%%%%%%%%%%%
%%%%%%%%%%%%%%%%%%%%%%%
%%%%%%%%%%%%%%%%%%%%%%%
%%%%%%%%%%%%%%%%%%%%%%%
%%%%%%%%%%%%%%%%%%%%%%%
%%%%%%%%%%%%%%%%%%%%%%%
%%%%%%%%%%%%%%%%%%%%%%%
%%%%%%%%%%%%%%%%%%%%%%%
%%%%%%%%%%%%%%%%%%%%%%%
%%%%%%%%%%%%%%%%%%%%%%%
%%%%%%%%%%%%%%%%%%%%%%%
%%%%%%%%%%%%%%%%%%%%%%%
%%%%%%%%%%%%%%%%%%%%%%%
%%%%%%%%%%%%%%%%%%%%%%%
%%%%%%%%%%%%%%%%%%%%%%%
%%%%%%%%%%%%%%%%%%%%%%%
%%%%%%%%%%%%%%%%%%%%%%%
%%%%%%%%%%%%%%%%%%%%%%%

%%%%%%%%%%%%%%%%%%%%%%%%%%%%%%%%%%%%%%%%%%%%%%%%%%%%%%%%%%%%%%%%%%%%%%%%

\begin{slide}{String Rewriting Problem}

\begin{itemize}
\item let $\cA$ be a finite alphabet

\medskip
\item let $\sim$ be a symmetric relation on substrings of length at most $k$ (where $k$ is a constant)

\medskip
\item let $A$ be the adjacency matrix of the undirected graph $G=(V,E)$, where
\begin{itemize}

\medskip
\item the set of vertices $V$ is $\cA^r$ (strings of length $r$ over $\cA$)

\medskip
\item the set of edges $E$ corresponds to possible conversions between strings, that is,
\[
(avb,awb)\in E \quad \Longleftrightarrow \quad v\sim w\
\]
\end{itemize}
\end{itemize}
\end{slide}

%%%%%%%%%%%%%%%%%%%%%%%%%%%%%%%%%%%%%%%%%%%%%%%%%%%%%%%%%%%%%%%%%%%%%%%%

\begin{slide}{String Rewriting Problem}
\begin{itemize}
\item let $s$, $t$, and $t'$ be three strings in $\cA^r$
\item let 
\[
\Delta_{s,t,t'}(n):=A_{s,t}^n - A_{s,t'}^n
\]
be the difference between the number of possibilities to obtain $t$ from $s$ and those to obtain $t'$ from $s$ after exactly $n$ replacements
\item given the promise
\begin{itemize}
\item $|\Delta_{s,t,t'}(n)| \le c^n$ for all $n\in\N$ ({\em growth condition})
\item $|\Delta_{s,t,t'}(m)|\ge \epsilon c^m$
\end{itemize}
determine the sign of $\Delta_{s,t,t'}(m)$ where

$c>0$ is a constant, $m\in\N$ with $m=poly(r)$, and $\epsilon\in [0,1]$ with $\epsilon=poly(1/r)$
\end{itemize}
\end{slide}

%%%%%%%%%%%%%%%%%%%%%%%%%%%%%%%%%%%%%%%%%%%%%%%%%%%%%%%%%%%%%%%%%%%%%%%%

\begin{slide}{Motivation from Bioinformatics}
\begin{itemize}
\item strings $s$, $t$, and $t'$ correspond to DNA sequences 
\item mutations change substrings (possible mutations are specified by the relation $\sim$)
\item String rewriting problem:

Are there more possibilities to mutate from $s$ to $t$ or from $s$ to $t'$ if exactly $m$ mutations occur?
\end{itemize}
\end{slide}

%%%%%%%%%%%%%%%%%%%%%%%%%%%%%%%%%%%%%%%%%%%%%%%%%%%%%%%%%%%%%%%%%%%%%%%%

\begin{slide}{PromiseBQP-complete}
Theorem: String Rewriting is PromiseBQP-complete

\bigskip
The string rewriting problem is the most difficult problem that can be efficiently solved on a quantum computer.

\bigskip
Quantum computers are better at estimating differences of combinatorial quantities than classical computers (provided that BPP is not equal to BQP).
\end{slide}

%%%%%%%%%%%%%%%%%%%%%%%%%%%%%%%%%%%%%%%%%%%%%%%%%%%%%%%%%%%%%%%%%%%%%%%%

\begin{slide}{String Rewriting is in PromiseBQP}

\begin{itemize}
\item rewrite $\Delta_{s,t,t'}(m)$ as a difference of expected values
\[
\Delta_{s,t,t'}(m)=\sqrt{2}(\<\psi^+|A^m|\psi^+\> - \<\psi^-|A^m|\psi^-\>)
\]
with $|\psi^{\pm}\>=(|s\>\pm|\varphi\>)/\sqrt{2}$ and $|\varphi\>=(|t\>+|t'\>))/\sqrt{2}$
\item determine the expected values as follows:
\begin{itemize}
\item let $U:=\exp(-i A/b)$ where $b\le\|A\|$
\item apply quantum phase estimation to $U$ and $|\psi^{\pm}\>$

\medskip
this allows one to sample according to the spectral measure induced by $A$ and $|\psi^{\pm}\>$: 

eigenvalue $\lambda_j$ occurs with probability  $\<\psi^{\pm}|P_j|\psi^{\pm}\>$
\end{itemize}
\item raise the outcome to the $m$th power
\item output average value after sufficiently many runs 
\end{itemize}
\end{slide}